\section{Abstract-Stage Triage}
\label{sec:ECTriage}

The venue sweep returns \(379\) records at the eleven venues. Most are not about contextual
optimization at all: the query is deliberately recall-oriented (\cref{sec:ECSearch}), and terms such
as \emph{task loss} and \emph{differentiable optimization} are common in unrelated machine-learning
work. Before obtaining full texts we therefore apply a single coarse filter, which we call
\emph{abstract-stage triage}.

We stress at the outset what this stage does \emph{not} do. It does not apply the inclusion rules of
\cref{sec:ECInclusion}, and it does not determine whether a paper studies a stochastic parameter. It
decides only which full texts are worth obtaining. Every substantive classification in this audit is
made from the full text, for reasons documented below.

\subsection{Why the Filter Is Not a Keyword Rule}

The natural implementation is a keyword or topic rule applied to the abstract, and we initially
adopted one: exclude a paper when its abstract indicates that no parameter is predicted, so that the
optimization instance is fully observed and the learning task is to solve it faster. Papers of this
kind -- amortized optimization, learned warm starts, neural combinatorial solvers -- are genuinely
out of scope, since with no unobserved \(Y\) there are no two oracles and no premium.

That rule fails systematically, and it fails on exactly the papers we can least afford to lose. We
tested it against twelve papers already known to be in scope, judging each from its abstract alone.
Four of the twelve would have been excluded: \citet{berthet2020learning},
\citet{poganvcic2020differentiation}, \citet{sahoo2023backpropagation} and
\citet{niepert2021implicit}. Each describes its
\emph{method} -- differentiating through a combinatorial solver -- and leaves the \emph{setting}, in
which a cost vector is predicted from context, to its experiments. Their abstracts contain no
language about prediction at all.

Abstracts at machine-learning venues foreground methodological contribution, so the papers most
likely to be misread are precisely those in the differentiable-optimization lineage that the search
was widened to capture. A referee has no way to
detect such an exclusion from the published record, since an excluded paper leaves no trace. We
therefore adopt a rule that never asks whether anything is predicted.

\subsection{The Triage Rule}

A record is retained when both of the following hold: an optimization or decision problem is the
object of the work, and a learned component appears in the same pipeline. A record is excluded when
there is no decision problem at all, when the only optimization is machinery for training or
compressing a model (hyperparameter search, architecture search, pruning, distillation), or when the
work is numerical simulation of a physical system. Records that cannot be judged are retained.

The asymmetry is deliberate. A record wrongly retained costs one full-text retrieval; a record
wrongly excluded is invisible for the remainder of the study.

\subsection{Implementation and Reproducibility}

The rule requires reading an abstract for what a paper is about, which no lexical pattern captures.
We implement it as a language-model classifier under a fixed prompt, and treat the prompt as part of
the method rather than as an implementation detail: it is versioned in the replication repository,
states the criterion as enumerated inclusions and exclusions rather than as an abstract principle,
and carries five worked examples with their reasoning, among them \citet{berthet2020learning} as a
retained record. Two instructions in the prompt exist solely to prevent the failure above, and each
records the evidence that motivated it.

The model identifier is pinned and decoding is greedy. For each record we store the assigned
bucket, a one-sentence rationale in the model's own words, a self-reported confidence, the model
identifier, a hash of the prompt, and the run date. The rationale is what makes the stage auditable:
a reader disagreeing with a particular exclusion can see the stated ground for it.

Because greedy decoding does not make a language model deterministic, every record is classified
three times and assigned its majority bucket, with the three individual buckets retained so that a
split is visible rather than silently resolved. Re-classifying one hundred records under the pinned
model returned the assigned bucket in all \(100\) cases, with the stated rationale identical in
\(76\). Note what this compares: a fresh single classification against the three-run majority of
record, which is the quantity that enters the labels we use. It concerns decoding variation alone,
and we are explicit below that it is not the quantity a reader should care most about.

We are careful not to overclaim here. Greedy decoding reduces run-to-run variation but does not
eliminate it, and a provider-side model revision can move a borderline record. This is why the stage
is validated on every execution against a labelled control set of papers known to be in scope,
four of which -- the four named above -- carry no prediction language in their abstracts and are the
binding test; the pipeline halts if any control record is excluded. We report the control-set result
alongside the triage counts.

\subsection{Execution and Outcome}

We executed the stage once on the \(379\) sweep records, under the pinned model identifier
\texttt{claude-opus-4-5-20251101}. The three runs agreed unanimously on \(373\) records. Three
records are proceedings front matter, removed by a prior rule rather than by the classifier. Three
records split, and were ruled by hand; in all three the hand ruling confirmed the majority bucket, so
they are recorded as adjudications rather than as corrections.

A fourth record was re-ruled. It had been examined during the design phase described below, under a
version of the prompt that was subsequently revised; the revision moved it from excluded to retained,
and we judged that a verdict rendered against a superseded specification could not be carried
forward unexamined. We re-read it under the current rule and excluded it. We record this because the
general point governs any audit of this kind: a human verdict is a verdict about a stated rule, and
it expires when the rule changes.

After adjudication the stage retains \(158\) records and excludes \(221\). Adjudications are stored
in a separate file and applied programmatically, so the classifier's output is never edited and every
human override is visible as a difference between two files. No record belonging to the survey reference list --
the references of the two surveys from which the sampling frame is built (\cref{sec:ECSearch}) -- was
excluded. This is a consistency check rather than a validation, since a survey's reference list is
not a gold standard for our inclusion rules, but a failure here would have been informative.

\subsection{Validating the Excluded Records}

Triage errors in the retaining direction correct themselves at full text. Errors in the excluding
direction do not, and they are therefore the quantity we measure. We measure them in two phases, and
we keep the phases separate.

\textbf{Design.} The first phase reports no statistic. We reviewed every excluded record whose
recorded confidence was below high, together with every excluded record at an operations-research
venue, where off-topic noise is rare: thirteen in all. Twelve confirmed the exclusion. The thirteenth
was a paper whose learned component is an amortized solver for optimal transport, excluded on the
stated ground that optimal transport there is a computational tool rather than a decision problem.
That is a failure to follow the rule rather than a defect in the rule, and it exposed an ambiguity in
how the rule was stated: the requirement of a learned component was being read as requiring a
predictor of problem \emph{data}, so a learned \emph{solver} did not register as satisfying it. We
revised the prompt to say that the learned component may be the solver itself, keeping both
requirements in force -- a faster classical algorithm with nothing learned remains excluded, as does
numerical simulation of a physical system with no decision problem.

That revision moved \(15\) of the \(379\) records, which deserves comment: a filter whose output
shifts under a small change to its instructions invites the suspicion that it is unstable. The prompt
\emph{is} the operational definition of the rule, so a revision that moved nothing would mean it had
not taken effect. The question is instead whether the records that moved are the ones the revision
was about. Comparing the runs before and after, all fifteen lie on the boundary the revision redrew:
six newly retained (amortized solvers for optimal transport and for mixed-integer programs, learned
molecular docking), six newly excluded (bilevel optimization used as training machinery,
treatment-effect estimation, learned navigation with no optimization problem formulated), and three
moving through the undecided bucket. None lie elsewhere. We read this as a targeted edit rather than
as drift.

\textbf{Measurement.} With the prompt frozen, we drew a fresh random sample of \(50\) of the \(221\)
excluded records under a recorded seed, and one author classified them by hand from title, venue,
year and abstract alone, with the model's bucket, rationale and confidence withheld so that the
reading was independent of them. The draw, the reading sheet and the recorded verdicts are in the
replication repository, so the sample is reproducible rather than merely described. All \(50\) were
confirmed excluded. Since the sample is a substantial fraction of a finite population, the
appropriate bound is hypergeometric rather than binomial: with \(95\%\) confidence, at most \(11\) of
the \(221\) excluded records are wrongly excluded. (The corresponding Clopper--Pearson bounds, which
discard the finite-population information, are \(5.8\%\) one-sided, or \(12.9\) records, and
\(7.1\%\) two-sided.)

The two phases cannot be pooled. Targeted sampling deliberately oversamples where errors are
expected, so its error rate does not estimate the rate in the excluded pile; and a sample that drove
a revision of the instrument is contaminated as a measurement of it.

\subsection{Scope of the Claim}

Triage is a retrieval decision and we claim no more for it. The bound above is a bound on
disagreement between the classifier and one reader applying the written rule. It is not evidence that
the rule is the right one, which is argued in \cref{sec:ECInclusion} rather than estimated, and it is
not a claim about language-model triage in general: it is a statement about this prompt, at the
recorded hash, under the recorded model. We have not measured sensitivity to paraphrases of the
prompt that preserve the rule. We did pilot the rule against a smaller model of the same family
before pinning the one above, and that pilot is worth reporting. Under the identical frozen prompt
the smaller model differed on \(23\) of the \(379\) records, \(17\) of them exclusions the pinned
model did not make, and its own stated grounds say what went wrong: in thirteen of the seventeen it
recorded that no optimization or decision problem was the object of the work, in papers where the
optimization plainly is the object -- Bayesian optimization, model predictive control, inverse
optimization, performative prediction, learned optimal transport. This is the failure the revision
described above was written against, reappearing under the revised wording that was meant to close
it: the requirement of a learned component was again being read as a requirement that problem
\emph{data} be predicted, so an optimization that is itself the object of the work registered as
machinery for something else. The smaller model passed the control set, \(12/12\), under the prompt
against which it was piloted, and we do not report that as evidence in its favour: a model can pass
a twelve-paper control set and still misread the rule at the boundary in the way just described.
The control set is a floor -- it catches a model that has stopped recognizing the papers this stage
exists to protect -- and not a test of whether a model applies the rule where the rule is hard. That
is why the identifier is pinned rather than merely validated, and why every count reported here is
tied to a fixed identifier and a fixed prompt hash. Perhaps the most honest statement is that
this stage is designed so that its failures are recoverable, not so that it is correct: every
subsequent stage of the audit works from full texts, and the screening log records a ground for
each exclusion so that a reader may disagree with any of them individually.
